\section{FINAL DISCUSSION}

\subsection{Interpretation of the Results}

Across all the simulated systems, Explicit Euler consistently showed the poorest conservation of both energy and
angular momentum, meanwhile Leapfrog, Velocity Verlet and Symplectic Euler demonstrated substantially greater
long-term stability.
Leapfrog-4 shows the smallest energy errors at smaller timesteps without behaviour deterioration at coarser resolutions.
This highlights the significance of higher-order symplectic methods. Although, with a computation time trade off.

Explicit Euler demonstates such behaviour due to the accumulated truncation errors over each timestep.
Symplectic methods by their mathematical design preserve the Hamiltonian geometry.
Instead of estimating a precise value of energy, symplectic methods revolve around the true value in the sinusoid.
For example, the RK4 algorithm minimises the local truncation error with the higher resolution, yet doesn't preserve
the phase space.
Leapfrog-4 contains its symplectic algorithmic nature and geometry, which already preserves phase-space by design
and with a higher resolution can produce more accurate results.
Hence why, there is a difference between local accuracy and long-term stability.

\subsection{Leapfrog-4 Behaviour}

The fourth-order Leapfrog method has both excellent short-term and long-term performance in the two-body simulations.
It's symplecticity can guarantee a long-term performance in a simulation, although with a big trade-off of computational cost.
As earlier results have showed, Leapfrog-4 is about 10 second slower than it's second-order variation.

\subsection{Computational Cost v. Accuracy}

Although RK4 previously produced the lowest numerical error at fine resolution, this came at approximately four force evaluations
per timestep and the highest runtime among all methods.
Leapfrog and Velocity Verlet therefore provide a considerably better compromise between computational efficiency and
conservation accuracy.
The higher-order Leapfrog-4 not only has showed a quicker computer time for the same set of problems succeeding RK4 by 20 second,
but also conserving the data more accurately long-term than RK4.

\subsection{Limitations}

While the results presented here are consistent with established theory on symplectic integration, several
limitations constrain the scope of the conclusions and point towards natural extensions of this work.
Some previous limitations have been extended in the study, such as adding a symplectic higher-order method.
Although, the study still lacks variable timestep integration, N-Body problem setups, computational scaling
and Newtonian approximation, which all were mentioned earlier.